%==============================================================
%  Chapter 2 -- Scalable K-Step Markov-Chain Recommendation
%                via Matrix Factorization
%  Source: Semester-1 mid-term deck (Matrix Factorization).
%==============================================================
\chapter{Scalable \texorpdfstring{$K$}{K}-Step Markov-Chain Recommendation via Matrix Factorization}
\label{ch:markov-factorization}

This chapter develops the first line of work undertaken for the IndiaMART
recommendation problem: modelling product-to-product navigation as a Markov
chain and using its $K$-step transition behaviour to generate
recommendations. The central difficulty is one of scale. The transition
matrix is of the order of tens of millions of products on a side, and the
naive computation of its $K$-step powers is infeasible both in time and in
memory. We motivate a matrix-factorization reformulation that expresses the
$K$-step transition matrix in terms of two thin factors, analyse the memory
it saves, and then examine three candidate factorization schemes ---
Alternating Least Squares (ALS), Non-negative Matrix Factorization (NMF), and
Weighted NMF --- comparing them in the light of the non-negativity and
sparsity structure that the problem imposes.

%--------------------------------------------------------------
\section{Problem Statement}
\label{sec:mf-problem}

The starting point is a product-to-product association matrix
$M \in \mathbb{R}^{N \times N}$, where each entry $M_{ij}$ encodes the
strength of the transition from product $i$ to product $j$ --- in our setting,
the (empirical) probability that a user who interacts with product $i$ next
interacts with product $j$. The goal is to build a \emph{$K$-step Markov-chain
recommender}: rather than recommending only the immediate successors of a
product, we follow the chain forward $K$ steps so that recommendations capture
multi-hop navigation patterns across the catalogue.

Two structural properties of $M$ dominate the design:

\begin{itemize}
  \item \textbf{Scale.} The matrix is approximately
        $2.3~\text{crore} \times 2.3~\text{crore}$, i.e.\ roughly
        $N \approx 2.3 \times 10^{7}$ on each axis.
  \item \textbf{Sparsity.} The density is approximately $10^{-4}$; the
        overwhelming majority of product pairs never co-occur in a session.
\end{itemize}

The combination is what makes the problem hard. While $M$ itself is sparse and
storable, its powers are not: as the chain is advanced through successive
steps, paths between products proliferate and the density of $M^{k}$ grows
\emph{exponentially} with $k$. Within a few steps the matrix becomes
effectively dense, and materialising it would require an amount of memory that
no commodity machine can provide --- in practice leading to buffer overflow
and out-of-memory failure. The remainder of the chapter is concerned with
computing the $K$-step transition behaviour \emph{without} ever forming a
dense $M^{k}$.

%--------------------------------------------------------------
\section{Background}
\label{sec:mf-background}

\subsection{Recommender Systems}
A recommender system is a data-driven system that surfaces items relevant to a
user by exploiting patterns and relationships latent in interaction data.
Broadly, three families are distinguished. \emph{Content-based} methods
recommend items similar to those a user has already engaged with, using item
attributes. \emph{Collaborative filtering} methods rely instead on the
behaviour of the user population, recommending items favoured by users with
similar histories. \emph{Hybrid} approaches combine the two to offset their
individual weaknesses. Recommender systems are now ubiquitous in e-commerce,
streaming, and social platforms. They are nonetheless beset by recurring
challenges: the \emph{cold-start} problem for new users or items, \emph{data
sparsity} when interactions are few relative to the catalogue, and
\emph{scalability} as catalogues grow to the scale considered here. The work
in this thesis sits squarely against the last two of these.

\subsection{Markov Chains and \texorpdfstring{$K$}{K}-Step Transitions}
\label{sec:markov}
A Markov chain is a stochastic model of a sequence of events in which the
probability of the next event depends only on the present state and not on the
history that preceded it (the \emph{memoryless}, or first-order Markov,
property). The dynamics are captured by a transition matrix $M$, whose entry
$M_{ij}$ gives the probability of moving from state $i$ to state $j$ in a
single step.

The key fact we exploit is how multi-step transitions compose. The
probability of moving from state $i$ to state $j$ in exactly $k$ steps is the
$(i,j)$ entry of the $k$-th matrix power:
\begin{equation}
  \label{eq:kstep}
  \bigl[M^{k}\bigr]_{ij}
    \;=\; \Pr\!\left(\text{state } j \text{ after } k \text{ steps}
                     \mid \text{state } i \text{ now}\right).
\end{equation}
Equation~\eqref{eq:kstep} is exactly the object a $K$-step recommender needs:
ranking the row $\bigl[M^{K}\bigr]_{i,\,\cdot}$ yields the products most likely
to be reached from product $i$ within $K$ steps. The computational obstacle of
Section~\ref{sec:mf-problem} is now precise --- we must read off rows of
$M^{K}$ without forming $M^{K}$ explicitly.

%--------------------------------------------------------------
\section{Proposed Solution: Factorized \texorpdfstring{$K$}{K}-Step Powers}
\label{sec:mf-solution}

Direct iteration of a $2.3\times10^{7}$ square matrix to the $K$-th power is
impossible at this scale. Our proposal is to replace $M$ by a low-rank
factorization and to push the matrix power onto the (tiny) factors.

Suppose $M$ admits a rank-$r$ factorization
\begin{equation}
  \label{eq:factorize}
  M \;\approx\; U V^{\top},
  \qquad U \in \mathbb{R}^{N \times r},\;\; V \in \mathbb{R}^{N \times r},
\end{equation}
with $r \ll N$ (in our experiments $r$ is of the order of $100$). The two
factors $U$ and $V$ are thin: their storage is linear in $N$ rather than
quadratic. The power of this representation is that the $K$-step matrix can be
expressed entirely through a small $r \times r$ core. Substituting
\eqref{eq:factorize} and using associativity,
\begin{align}
  M^{k}
    &\approx \underbrace{(U V^{\top})(U V^{\top}) \cdots (U V^{\top})}_{k\ \text{factors}}
     \nonumber\\[2pt]
    &= U \,\bigl(V^{\top} U\bigr)^{k-1}\, V^{\top}
     \;=\; U\, W^{\,k-1}\, V^{\top},
  \label{eq:factored-power}
\end{align}
where we have defined the \emph{core matrix}
\begin{equation}
  \label{eq:core}
  W \;=\; V^{\top} U \;\in\; \mathbb{R}^{r \times r}.
\end{equation}
Equation~\eqref{eq:factored-power} is the heart of the approach: the only
power that must be computed is $W^{\,k-1}$, an $r\times r$ matrix with
$r \approx 100$, which is trivial. The large factors $U$ and $V$ are touched
only by the final, lazy left- and right-multiplications, and a single
recommendation row $\bigl[M^{K}\bigr]_{i,\cdot}
= U_{i,\cdot}\,W^{\,K-1}\,V^{\top}$ can be evaluated on demand without
materialising the full $N \times N$ result.

\begin{remark}
The mid-term presentation wrote the relation compactly as
$M^{k} = U\,(W^{k})\,V^{\top}$. The exact algebra in
\eqref{eq:factored-power} gives the exponent $k-1$ rather than $k$, since one
factor of $V^{\top}U$ is consumed by the outer $U$ and $V^{\top}$. The
distinction does not affect the memory argument below.
\end{remark}

\subsection{Memory Analysis}
\label{sec:mf-memory}
The saving is best appreciated with the concrete figures of the IndiaMART
setting. Take $N = 2.3 \times 10^{7}$, rank $r = 100$, step count $K = 6$, and
double-precision storage at $8$ bytes per entry.

If $M^{k}$ were materialised densely --- which, as noted, it effectively
becomes after a few steps --- it would occupy
\begin{equation}
  N^{2} \times 8
    = \bigl(2.3\times10^{7}\bigr)^{2}\times 8
    \;\approx\; 4.232 \times 10^{15}\ \text{bytes},
\end{equation}
i.e.\ on the order of petabytes, which is plainly infeasible. In the
factorized form, by contrast, we store only $U$, $V$, and the core $W$:
\begin{equation}
  \underbrace{2\,(N \times r)\times 8}_{U,\,V}
   + \underbrace{r^{2}\times 8}_{W}
    \;=\; 2\,(2.3\times10^{7}\times 100)\times 8 + \text{(negligible)}
    \;\approx\; 36.8\ \text{GB}.
\end{equation}
The reduction is from petabytes to tens of gigabytes --- five orders of
magnitude --- and brings the $K$-step computation within reach of a single
large-memory machine. Table~\ref{tab:mf-memory} summarises the comparison.

\begin{table}[ht]
  \centering
  \caption{Memory footprint of the $K$-step transition computation
           ($N=2.3\times10^{7}$, $r=100$, double precision).}
  \label{tab:mf-memory}
  \begin{tabular}{lll}
    \toprule
    Representation & Quantity stored & Memory \\
    \midrule
    Dense $M^{k}$            & $N^{2}$ entries          & $\approx 4.232\times10^{15}$ B \\
    Factorized $U W^{k-1}V^{\top}$ & $2Nr + r^{2}$ entries & $\approx 36.8$ GB \\
    \bottomrule
  \end{tabular}
\end{table}

%--------------------------------------------------------------
\section{Factorization Methods}
\label{sec:mf-methods}

The reformulation of Section~\ref{sec:mf-solution} presumes a factorization
$M \approx U V^{\top}$. How that factorization is obtained is itself a design
choice, and three methods were considered: Alternating Least Squares (ALS),
Non-negative Matrix Factorization (NMF), and Weighted NMF. Because the entries
of $M$ are transition probabilities, they are non-negative by construction ---
a property that, as we shall see, decisively favours the non-negative methods.

\subsection{Alternating Least Squares (ALS)}
\label{sec:als}
ALS is a classical matrix-factorization technique, widely used in
collaborative filtering and recommender systems~\cite{HKV08,KBV09}. It
decomposes a large, sparse matrix into lower-dimensional latent factors by
minimising the regularised reconstruction error
\begin{equation}
  \label{eq:als-obj}
  \min_{U,\,V}\;
    \bigl\| M - U V^{\top} \bigr\|_{F}^{2}
    \;+\; \lambda\bigl(\|U\|_{F}^{2} + \|V\|_{F}^{2}\bigr),
\end{equation}
where $\lambda > 0$ is a regularisation weight and $\|\cdot\|_{F}$ the
Frobenius norm. The objective in \eqref{eq:als-obj} is not jointly convex in
$(U,V)$, but it \emph{is} convex in each factor with the other held fixed.
ALS exploits this by alternating: it fixes $V$ and solves a least-squares
problem for $U$, then fixes $U$ and solves for $V$, repeating until
convergence. Each subproblem has a closed form,
\begin{align}
  \label{eq:als-update}
  u_{i} &= \bigl(V^{\top}V + \lambda I_{r}\bigr)^{-1} V^{\top} m_{i},
   &\text{(row $i$ of $U$)} \nonumber\\
  v_{j} &= \bigl(U^{\top}U + \lambda I_{r}\bigr)^{-1} U^{\top} m_{j},
   &\text{(row $j$ of $V$)}
\end{align}
where $m_{i}$ and $m_{j}$ denote the relevant row/column of $M$ and $I_{r}$ is
the $r\times r$ identity. ALS is efficient, scalable, parallelises naturally
across rows, and copes well with sparse data, which is why it is a standard
first choice for recommendation at scale.

\paragraph{Why ALS does not work here.}
Despite its appeal, ALS proved unsuitable for this problem. The entries of $M$
are transition \emph{probabilities} and are therefore non-negative by
definition; a faithful factorization should respect this. ALS, however, places
no sign constraint on $U$ or $V$, and the reconstructed $UV^{\top}$ readily
produces negative entries that have no probabilistic meaning. In the
experiments conducted, approximately $34\%$ of the reconstructed entries were
negative, and the normalised reconstruction error was $51.23\%$. Both figures
are unacceptable for a probability matrix. This observation is what motivates
moving to non-negativity-constrained factorizations.

\subsection{Non-Negative Matrix Factorization (NMF)}
\label{sec:nmf}
NMF decomposes a non-negative matrix $M$ into two lower-dimensional factors
that are \emph{themselves} constrained to be non-negative,
$M \approx U V^{\top}$ with $U, V \ge 0$ elementwise. The objective is
\begin{equation}
  \label{eq:nmf-obj}
  \min_{U,\,V \,\ge\, 0}\;
    \bigl\| M - U V^{\top} \bigr\|_{F}^{2}.
\end{equation}
The non-negativity constraint guarantees that the reconstruction
$U V^{\top}$ is itself non-negative, so every entry remains interpretable as a
(non-normalised) probability --- precisely the property ALS failed to
preserve. The standard solver is the multiplicative-update rule of Lee and
Seung~\cite{LeeSeung01}, which preserves non-negativity automatically at every
iteration:
\begin{equation}
  \label{eq:nmf-update}
  U \;\leftarrow\; U \;\odot\; \frac{M V}{\,U V^{\top} V\,},
  \qquad
  V \;\leftarrow\; V \;\odot\; \frac{M^{\top} U}{\,V U^{\top} U\,},
\end{equation}
where $\odot$ denotes the Hadamard (elementwise) product and the division is
elementwise. Starting from non-negative initial factors, the updates in
\eqref{eq:nmf-update} keep $U$ and $V$ non-negative throughout and
monotonically decrease the objective \eqref{eq:nmf-obj}.

\subsection{Weighted NMF (WNMF)}
\label{sec:wnmf}
Weighted NMF extends NMF by attaching a weight to each entry of $M$, so that
different observations contribute with different importance. This is valuable
when the data contains missing values, noise, or entries of varying
confidence --- all of which characterise an extremely sparse association
matrix. Let $W^{\mathrm{wt}} \in \mathbb{R}^{N\times N}_{\ge 0}$ be the weight
matrix. The objective becomes
\begin{equation}
  \label{eq:wnmf-obj}
  \min_{U,\,V \,\ge\, 0}\;
    \bigl\| W^{\mathrm{wt}} \odot \bigl(M - U V^{\top}\bigr) \bigr\|_{F}^{2},
\end{equation}
and the corresponding multiplicative updates are
\begin{equation}
  \label{eq:wnmf-update}
  U \;\leftarrow\; U \;\odot\;
    \frac{\bigl(W^{\mathrm{wt}}\odot M\bigr) V}
         {\bigl(W^{\mathrm{wt}}\odot (U V^{\top})\bigr) V},
  \qquad
  V \;\leftarrow\; V \;\odot\;
    \frac{\bigl(W^{\mathrm{wt}}\odot M\bigr)^{\top} U}
         {\bigl(W^{\mathrm{wt}}\odot (U V^{\top})\bigr)^{\top} U}.
\end{equation}
Setting $W^{\mathrm{wt}} = \mathbf{1}$ (all weights equal) recovers ordinary
NMF; the freedom lies in choosing $W^{\mathrm{wt}}$ to reflect the structure of
the data.

%--------------------------------------------------------------
\section{Results and Discussion}
\label{sec:mf-results}

The decisive design lever turned out to be the choice of weights in WNMF. By
taking \emph{binary} weights --- $W^{\mathrm{wt}}_{ij} = 1$ where $M_{ij}$ is
observed and $0$ where it is missing --- the factorization is asked to fit only
the entries that actually carry information, and the vast field of structural
zeros in the sparse matrix is prevented from dominating the loss. This removes
the undesired influence of missing elements on the factors. Under binary
weighting the reconstruction error on the observed entries dropped almost to
zero, a dramatic improvement over both ALS (Section~\ref{sec:als}) and
unweighted NMF.

% FIGURE PLACEHOLDER: user to supply the error-comparison plot (ALS vs NMF
% vs binary-weighted WNMF) into figures/.
% \begin{figure}[ht]
%   \centering
%   \includegraphics[width=0.8\textwidth]{figures/mf_error_comparison.png}
%   \caption{Normalised reconstruction error across the three factorization
%            methods.}
%   \label{fig:mf-error}
% \end{figure}

Table~\ref{tab:mf-comparison} collects the comparison qualitatively; the
quantitative error figures from the experiments are reported where available.

\begin{table}[ht]
  \centering
  \caption{Comparison of the three factorization methods for the non-negative,
           highly sparse transition matrix.}
  \label{tab:mf-comparison}
  \begin{tabular}{lccc}
    \toprule
    Method & Non-negativity & Handles missing entries & Reconstruction error \\
    \midrule
    ALS              & No (\,$\approx34\%$ negative\,) & No   & $51.23\%$ \\
    NMF              & Yes                              & No   & --- \\
    WNMF (binary $W$) & Yes                             & Yes  & $\approx 0$ \\
    \bottomrule
  \end{tabular}
\end{table}

\noindent The progression across the three methods tells a clear story. ALS,
the conventional choice, breaks on the non-negativity requirement intrinsic to
a probability matrix. NMF restores non-negativity but still tries to fit the
structural zeros of a $10^{-4}$-dense matrix. Binary-weighted NMF addresses
both issues at once and drives the error on the observed entries almost to
zero, which is why it was adopted as the working method at the end of the
mid-term stage. The closer evaluation carried out subsequently, however,
revealed that this near-zero error figure conceals a serious failure mode. The
next section examines it.

%--------------------------------------------------------------
\section{Limitation of Weighted NMF}
\label{sec:wnmf-limit}

The very mechanism that makes binary-weighted NMF attractive --- fitting only
the observed entries --- is also the source of its weakness. Recall from
\eqref{eq:wnmf-obj} that with a binary weight matrix the loss is accumulated
\emph{only} over positions where $W^{\mathrm{wt}}_{ij}=1$. The multiplicative
updates of \eqref{eq:wnmf-update} therefore move $U$ and $V$ so as to improve
the reconstruction of those observed positions and nothing else. Likewise, the
absolute (and normalised) error is measured only against the observed entries
of the original matrix. The reconstructed values at the \emph{unobserved}
positions --- those with $W^{\mathrm{wt}}_{ij}=0$ --- enter neither the loss nor
the error metric. They are produced as an uncontrolled by-product of the
factors and are never monitored during training.

This is precisely where the method fails. Because nothing in the objective
constrains the unobserved entries, there is no force keeping them small or
well-behaved. As the factors are driven to fit the observed entries, some of
the unconstrained entries in the reconstruction $U V^{\top}$ drift to
unusually large values. Two consequences follow:

\begin{itemize}
  \item \textbf{The error metric is blind to the problem.} Since error is
        computed only on observed entries, a model can report near-zero error
        while simultaneously placing wildly inflated values at unobserved
        positions. The optimistic figure of Section~\ref{sec:mf-results} is, in
        this light, partly an artefact of \emph{what} is being measured.
  \item \textbf{Inflated entries dominate the Markov chain.} In the $K$-step
        recommender these spuriously large reconstructed values behave as if
        they were high transition probabilities. Propagated through
        $M^{K}\approx U\,W^{\,K-1}\,V^{\top}$, they dominate the chain's
        transitions, so that a handful of products surface excessively in the
        recommendations regardless of the query --- crowding out the genuinely
        relevant items.
\end{itemize}

In short, binary-weighted NMF optimises a quantity (observed-entry
reconstruction) that is not aligned with the quantity that actually matters for
recommendation (the behaviour of the \emph{full} reconstructed matrix under
repeated multiplication). Controlling the unobserved entries --- so that the
reconstruction remains a well-behaved transition operator everywhere, not only
where data was seen --- is the open problem that this work carries into the
next stage.
